
\documentclass[letterpaper,11pt]{article}
%%%%%%%%%%%%%%%%%%%%%%%%%%%%%%%%%%%%%%%%%%%%%%%%%%%%%%%%%%%%%%%%%%%%%%%%%%%%%%%%%%%%%%%%%%%%%%%%%%%%%%%%%%%%%%%%%%%%%%%%%%%%%%%%%%%%%%%%%%%%%%%%%%%%%%%%%%%%%%%%%%%%%%%%%%%%%%%%%%%%%%%%%%%%%%%%%%%%%%%%%%%%%%%%%%%%%%%%%%%%%%%%%%%%%%%%%%%%%%%%%%%%%%%%%%%%
\usepackage{graphicx}
\usepackage{amsmath,amsfonts,amssymb,amsthm,float}
\usepackage{hyperref}
\usepackage[utf8]{inputenc}
\usepackage[left=2cm, right=2cm, top=2cm, bottom=2cm]{geometry}

\setcounter{MaxMatrixCols}{10}
%TCIDATA{OutputFilter=LATEX.DLL}
%TCIDATA{Version=5.50.0.2953}
%TCIDATA{<META NAME="SaveForMode" CONTENT="1">}
%TCIDATA{BibliographyScheme=BibTeX}
%TCIDATA{LastRevised=Saturday, February 21, 2026 23:39:12}
%TCIDATA{<META NAME="GraphicsSave" CONTENT="32">}
%TCIDATA{Language=American English}

\input{tcilatex}
\renewcommand{\baselinestretch}{1.15}
\setlength{\parindent}{0pt}
\setlength{\parskip}{0.5\baselineskip}
\pretolerance=2000 \tolerance=3000
\renewcommand{\abstractname}{Resumen}

\begin{document}

\title{Sistema presa-depredador de Lotka-Volterra}
\author{Lopez Pardo Brianna Vanessa$\left[ 22212261\right] $ \\
%EndAName
Departamento de Ingenier\'{\i}a El\'{e}ctrica y Electr\'{o}nica\\
Tecnol\'{o}gico Nacional de M\'{e}xico / Instituto Tecnol\'{o}gico de Tijuana%
}
\maketitle

\section{\noindent \textbf{Palabras clave: }Sistema din\'{a}mico;
Semitrayectoria; Lema de positividad; Invariante en el tiempo; EDO de primer
orden.}

\noindent El autor debe identificar palabras clave espec\'{\i}ficas que
representen con precisi\'{o}n la investigaci\'{o}n, se debe evitar usar t%
\'{e}rminos generales o vagos que podr\'{\i}an aplicarse a muchos temas
diferentes; se sugiere escribir de cuatro a siete palabras clave ordenadas
alfab\'{e}ticamente, separadas por punto y coma, y omitir palabras que ya est%
\'{a}n en el t\'{\i}tulo del documento.

\bigskip

\noindent Correo: \textbf{l22212261@tectijuana.edu.mx}

\noindent \noindent Carrera: \textbf{Ingenier\'{\i}a Biom\'{e}dica }

\noindent Asignatura: \textbf{Biolog\'{\i}a de Sistemas}

\noindent Profesor: \href{https://biomath.xyz/}{\textbf{Dr. Paul Antonio
Valle Trujillo}} (paul.valle@tectijuana.edu.mx)

\section{Modelo matem\'{a}tico}

El sistema presa-depredador de Lotka-Volterrra se formula mediante las
siguientes dos Ecuaciones Diferenciales Ordinarias (EDOS) no lineales de
primer orden:

\begin{eqnarray*}
\dot{x} &=&\alpha x-\beta xy, \\
\dot{y} &=&\delta xy-yx,
\end{eqnarray*}

donde los par\'{a}metros $\alpha ,\beta ,\delta ,\gamma >0$ y las
condiciones iniciales x$(0),y(0)\geq 0$. La variables $x(t)$ representa a
las presas [\textit{Hares}], la variables $y(t)$ representa a los
depredadores [\textit{Lynx}] y el tiempo $\ t$ se mide en a\~{n}os.

\bigskip

\section{An\'{a}lisis de positividad}

El an\'{a}lisis de positividad permite obtener conclusiones sobre la din\'{a}%
mica de las semitrayectorias positivas $\left( \Gamma ^{+}\right) $ [Secci%
\'{o}n 2.3] para condiciones iniciales negativas $\left[ x\left( 0\right)
,y\left( 0\right) \geq 0\right] ,$ para esto se debe evaluar al sistema en
la frontera del primer cuadrante (ortante positivo) como se indica a
continuaci\'{o}n:\ 

\begin{eqnarray*}
\left. \dot{x}\right\vert _{x=0} &=&\alpha (0)-\beta (0)y=0, \\
\left. \dot{y}\right\vert _{y=0} &=&\alpha (0)-\beta (0)y=0,
\end{eqnarray*}
\ \ \ \ \ \ \ \ \ \ \ \ \ por lo tanto, con base en el \textbf{Lema de
Positividad }en sistemas no lineales [Secci\'{o}n 2.6], se establece que el
sistema presa-depredador de Lotka-Volterra tiene soluciones positivamente
invariantes[Secci\'{o}n 3.1] para condiciones iniciales no negativas.

\section{Puntos de equilibrio}

Para calcular los puntos de equilibrio es necesario igualar cada ecuacion
del sistema a cero y resolver para cada una de las variables, es decir,

\begin{eqnarray*}
\alpha x-\beta xy &=&0 \\
\delta xy-\gamma y &=&0
\end{eqnarray*}

$\func{assume}(\alpha ,\func{positive})=\allowbreak \left( 0,\infty \right) $

$\func{assume}(\beta ,\func{positive})=\allowbreak \left( 0,\infty \right) $

$\func{assume}(\delta ,\func{positive})=\allowbreak \left( 0,\infty \right) $

$\func{assume}(\gamma ,\func{positive})=\allowbreak \left( 0,\infty \right) $

lo anterior permite que los cuatro par\'{a}metros del sistema son positivos.
Al aplicar el comando compute-solve-exact, se obtienen los siguientes puntos
de equilibrio:\ 

\bigskip 
\begin{eqnarray*}
(x_{1}^{\ast },y_{1}^{\ast }) &=&(0,0), \\
(x_{1}^{\ast },y_{1}^{\ast }) &=&\left( \frac{\gamma }{\delta },\frac{\alpha 
}{\beta }\right) , \\
&&\left[ x=\frac{\gamma }{\delta },y=\frac{\alpha }{\beta }\right] ,\left[
x=0,y=0\right]
\end{eqnarray*}

Nota: En algunos libros cuandoel origen del espaciode estados es un
equilibrio del sistema, se escibre as\'{\i}:

\begin{equation*}
(x_{0}^{\ast },y_{0}^{\ast })=(0,0)
\end{equation*}

\section{Estabilidad asintotica}

Para analizar la estabilidad local de los puntos de equilibrio de un sistema
din\'{a}mico [secciones 5.1 y 5.2.3.] es necesario calcular la matriz
Jacobiana [secci\'{o}n 2.7], como se inidica a continuaci\'{o}n:

\begin{equation*}
J=\left[ 
\begin{array}{cc}
\frac{\delta \dot{x}}{\delta x} & \frac{\delta \dot{x}}{\delta y} \\ 
\frac{\delta \dot{y}}{dx} & \frac{\delta \dot{y}}{\delta y}%
\end{array}%
\right] =\left[ 
\begin{array}{cc}
\alpha -\beta y & -\beta x \\ 
\delta y & \delta x-y%
\end{array}%
\right] .
\end{equation*}

Ahora, se evalua cada no de los puntos de equilibrio calculados, primero se
analiza $(x_{1}^{\ast }y_{1}^{\ast })=(0,0)$, obteniendo el siguiente
resultado:

\begin{equation*}
\left. J\right\vert _{(x_{1}^{\ast }y_{1}^{\ast })}=\ \left[ 
\begin{array}{cc}
\alpha -\beta y & -\beta x \\ 
\delta y & \delta x-y%
\end{array}%
\right] _{x=0,y=0}=\left[ 
\begin{array}{cc}
\alpha  & 0 \\ 
0 & -\gamma 
\end{array}%
\right] ,
\end{equation*}

entonces, debido a que el resultado es una matriz diagonal [secci\'{o}n
2.7], los valores propios est\'{a}n dados por cada uno de los elementos de
dicha diagonal , es decir, 

\begin{eqnarray*}
\lambda _{1} &=&\alpha , \\
\lambda _{2} &=&-\gamma ,
\end{eqnarray*}

con base en estos resultados se concluye lo siguiente: El punto de
equilibrio $(x_{1}^{\ast },y_{1}^{\ast })=(0,0)$ es inestable [secciones 5.1
y 5.2.3].

Al evaluar el segundo punto de equilibrio $(x_{2}^{\ast },y_{2}^{\ast })=(%
\frac{\gamma }{\delta },\frac{\alpha }{\beta })$, se obtiene el siguiente
resultado:   

\begin{equation*}
\left. J\right\vert _{(x_{2}^{\ast }y_{2}^{\ast })}=\ \left[ 
\begin{array}{cc}
\alpha -\beta y & -\beta x \\ 
\delta y & \delta x-y%
\end{array}%
\right] _{x=\frac{\gamma }{\delta },y=\frac{\alpha }{\beta }}=\left[ 
\begin{array}{cc}
\alpha -\beta \frac{\alpha }{\beta } & -\beta \frac{\gamma }{\delta } \\ 
\delta \frac{\alpha }{\beta } & \delta \frac{\gamma }{\delta }-\gamma 
\end{array}%
\right] =\left[ 
\begin{array}{cc}
0 & -\frac{\beta \gamma }{\delta } \\ 
\frac{\alpha \delta }{\beta } & 0%
\end{array}%
\right] ,
\end{equation*}

por lo tanto, se deben de calcular los valores propios mediante la siguiente
ecuaci\'{o}n [secci\'{o}n 2.7]:

\begin{equation*}
\det (A-\lambda I)=0,
\end{equation*}

donde $A$ es la matriz Jacobiana evaluada en el punto de equilibrio, $%
\lambda \in \mathbf{R}$ (es un coeficiente) e $\ I~$es la amtriz identidad,
al sustituir se obtiene lo siguiente:\ 

\begin{equation*}
\det \left[ \left( 
\begin{array}{cc}
0 & -\frac{\beta \gamma }{\delta } \\ 
\frac{\alpha \delta }{\beta } & 0%
\end{array}%
\right) -\lambda \left( 
\begin{array}{cc}
1 & 0 \\ 
0 & 1%
\end{array}%
\right) \right] =0,
\end{equation*}

entonces, al realizar las operaciones aritm\'{e}ticas necesarias se obtiene
el siguiente resultado:

\begin{eqnarray*}
\det \left[ \left( 
\begin{array}{cc}
0 & -\frac{\beta \gamma }{\delta } \\ 
\frac{\alpha \delta }{\beta } & 0%
\end{array}%
\right) -\lambda \left( 
\begin{array}{cc}
1 & 0 \\ 
0 & 1%
\end{array}%
\right) \right]  &=&0, \\
\det \left( 
\begin{array}{cc}
-\lambda  & -\frac{\beta \gamma }{\delta } \\ 
\frac{\alpha \delta }{\beta } & -\lambda 
\end{array}%
\right)  &=&0, \\
\lambda ^{2}+\alpha \gamma  &=&0,
\end{eqnarray*}

y los valores propios est\'{a}n dados por:

\begin{eqnarray*}
\lambda _{1} &=&0+\sqrt{-\alpha \gamma }=+i\sqrt{\alpha \gamma }, \\
\lambda _{2} &=&0-\sqrt{-\alpha \gamma }=-i\sqrt{\alpha \gamma },
\end{eqnarray*}

\bigskip 

\end{document}
